\subsubsection{Continuidad, recarga y fuentes por estado (CEN-4)}
\label{sec:sd4-continuidad-cen4}

La configuración mantiene dos Eaton \texttt{93T60KMBSB}, una por camino, y ochenta CSB \texttt{HRL12540W} por UPS, en dos cadenas de cuarenta con punto central. El suministro de todo el parque comprende TI, comunicaciones, captura, protección, dieciséis puestos y sus pantallas, motores, cargas móviles y los auxiliares ambientales que permiten conservar la autonomía. Los siete puestos prioritarios no reducen esa frontera. Cada camino sostiene por separado el parque al retirar el otro; los sesenta minutos iniciales y treinta minutos mínimos mantenidos se verifican con generación inhibida y la carga completa, no sumando bancos A/B \parencites[RT-06.07, p.~15]{pucv-transversal}{p7-eaton-93t-guia}{csbHRL12540WRA260131}.

\textbf{Decisión de ventilación y recarga.} La seguridad de gasificación se dimensiona para el máximo configurable publicado de 60 A agregados por UPS, con 900 m$^3$/h efectivos por banco. La propuesta deja de depender de que el servicio autorizado pueda fijar exactamente 20 A. Los 60 A son un techo agregado de la UPS y tampoco autorizan 60 A en una sola cadena CSB: el límite publicado es 54 A por bloque. Los 60 A no son una corriente simultáneamente disponible a cualquier carga: la entrada del rectificador, la temperatura y la capacidad de fuente siguen limitando la recuperación. La ventilación y su alimentación son permanentes en red, batería y grupo; no se admite apagarla al inhibir la carga porque el fabricante advierte liberación de gas también durante descarga. Su selección y tratamiento se desarrollan en ENV16. La autonomía de los ventiladores se recibe junto con el banco, sin atribuirle la autonomía propia del panel de incendio \parencites[p.~98 y advertencias de seguridad]{p7-eaton-93t-guia}{csbManualJunio2025}.

La guía permite encender y apagar el cargador desde el control local. Esto ofrece una intervención manual definida; no acredita un contacto automático dedicado de inhibición. REPO aísla la salida crítica y exige circuito dedicado: no se usa como señal de ventilación ni como sustituto de una inhibición selectiva del cargador. El proyecto requiere un enclavamiento local de dos niveles: aviso de caudal insuficiente o temperatura fuera de banda y retirada selectiva de recarga; persistencia de la condición requiere transferencia del servicio a una fuente compatible y aislamiento del banco afectado por personal competente. La disponibilidad del camino sano se comprueba antes de transferir. Con fallas simultáneas se mantiene la protección de personas y se declara degradación; no se promete continuidad indefinida desde un banco inseguro. La interfaz automática de inhibición y su aceptación por Eaton permanecen por individualizar; el comando manual documentado no se presenta como esa interfaz \parencite[secc.~6.2 y 6.5]{p7-eaton-93t-guia}.

\textbf{Dos máximos que no se suman como capacidades disponibles.} A 400 V, la cota de entrada publicada es $S_{in}=\sqrt{3}(400)(108)/1000=74,8246$ kVA. Con techo DC de proyecto de 588 V, 60 A representan 35,28 kW entregados al banco. Bajo eficiencia conjunta propia de 0,94, sin agregar auxiliares previos, 60 kW de salida más esa recarga exigirían $(60+35,28)/0,94=101,3617$ kW de entrada, superiores incluso a la cota aparente. La potencia nominal de salida y el máximo de cargador no describen un único estado realizable \parencites[pp.~33--35 y 98]{p7-eaton-93t-guia}{p7-eaton-93t-tecnica}.

Para selección, con $P_{in}\leq S_{in}$ y conservando la eficiencia propia, el límite superior de cribado es
\[
 I_{c,cribado}(P_{sal})=\max\left\{0,\frac{0,94(74,824595)-P_{sal}}{0,588}\right\}\ \mathrm{A}.
\]
Resultan 17,5767 A a 60 kW de salida y 37,1488 A a 48,4916 kW. Si se presume además FP de entrada 0,99, se obtienen 16,3806 y 35,9526 A. Son techos aritméticos con supuestos de proyecto, no curvas Eaton, ajustes disponibles ni garantía de corriente mínima. A 60 A, el mismo presupuesto permite como máximo 35,0551 kW de salida, o 34,3518 kW con FP 0,99. El servicio autorizado debe establecer el reparto real, las pérdidas y su límite de entrada; la recepción captura salida, entrada y DC simultáneos. Un tiempo de recarga calculado desde 60 A constantes a plena salida sería inválido.

\textbf{Mapa de alimentación y de fallas.} Servidores, NGFW y core mantienen PSU A/B. Los puestos seleccionados en PGE-4 conservan el cruce AC/DC de ambas pantallas y fuentes DC independientes del PC; no se añade otra partida de 370 W por duplicar fuentes. El panel INC-4 dispone de alimentación AC con transferencia A/B y reserva DC propia; se conserva servicio durante la maniobra sin exigir que la transferencia sea de tiempo cero. Los monitores ambientales MON16 se alimentan directamente desde A/L2 y B/L3, respectivamente, con dos redes completas de sondas que cubren todas las zonas; al retirar una fuente se conserva la otra red de monitorización. No se agrega una transferencia ni se atribuye doble alimentación a cada monitor. La seguridad de caudal y gas no depende de que esos monitores o la WAN estén disponibles. Las alimentaciones se reciben por pérdida de entrada, salida y fuente interna, que son fallas distintas.

Para NVR, terminaciones WAN, radio, motores, cargas móviles y conversiones de campo se conserva la partida individual hasta seleccionar entrada y sostenimiento. Una transferencia AC exige comparar su peor interrupción con el sostenimiento mínimo documentado del activo; para DC se exige doble entrada real o fuentes desacopladas autorizadas. Un STS genérico, dos cables conectados a la misma PSU o una batería de panel que no alimenta el resto del parque no cierran RT-06.07. El mapa registra tensión, frecuencia, fase, origen, protección, potencia por estado y conducta tras fallo; no descarta una carga imprescindible para recuperar margen eléctrico.

\subsubsection{Bancada seleccionada y reserva de generación (GCI-4)}
\label{sec:sd4-generacion-gci4}

Se seleccionan dos Avtron \texttt{4100}, una por grupo, con 60 kW resistivos a 400 V/50 Hz, resolución de 5 kW y opción Automatic Load Control (ALC); se reserva una tercera unidad compatible apagada. La ficha publica rango estándar de 10--150 kW, variante 400 V/50 Hz, servicio exterior continuo entre $-28$ y 50\,$^{\circ}$C, enclavamientos de presión/temperatura y tolerancia resistiva de $-0/+5\,\%$. El conjunto mide aproximadamente $1.219\times685\times889$ mm y pesa 227 kg; mantiene 2,44 m libres de obstáculos en admisión/descarga y solera exterior. Su calor se descarga fuera de TI, bancos, admisión de grupos y rutas de combustible. La capacidad de 60 kW no significa conectar todos los escalones junto con 90 kW de servicio \parencite[pp.~1--2]{lote17-avtron-4100}.

El ALC incorpora CT sobre la carga de servicio aguas abajo de la derivación de bancada; ese CT no mide el total del grupo. Inicia por contacto normalmente abierto y añade/retira escalones con retardos. Su ventana depende del mayor escalón utilizado, que no se deduce de una resolución de 5 kW. Se añade medición del total del grupo y permiso local de escalón que compruebe kW, kVA, corriente de fase y margen transitorio; retirar la bancada tiene prioridad frente a arranque climático o recuperación de recarga. La consigna propia inicial de 57,6 kW no se atribuye a FG Wilson ni al 60\,\% típico del ejemplo ALC. El fabricante del grupo confirma umbral, duración y mantenimiento por baja carga. La bancada aporta carga activa; no reproduce por sí sola armónicos, excitación reactiva ni toda la respuesta dinámica \parencite[pp.~1--2]{lote17-avtron-alc}.

La variante resistiva 400 V/50 Hz está publicada, pero la ficha describe auxiliar de control 120 V/60 Hz y ventilación sin potencia de motor individualizada. No se conecta ese auxiliar directamente a 230 V/50 Hz ni se presupone que un transformador cambie frecuencia. La liberación de la configuración 4100 exige diagrama específico de motor/control para 50 Hz, potencia/corriente de auxiliares, escalones reales y CT ajustado a cada P200-6. La opción de transformador, por sí sola, no resuelve la frecuencia. Esta frontera queda abierta y se identifica separadamente de un ensayo futuro; no se presenta una bancada seleccionada como conjunto auxiliar ya compatible.

Se conservan límites totales de 90 kW/90 kVA por grupo y requisito de 129 kW efectivos a la admisión de proyecto. A 400 V, 90 kVA representan 129,9038 A, por debajo de la derivación limitada a 160 A. La bancada máxima de 60 kW supone 86,6025 A antes de auxiliares y 63 kW en su tolerancia superior; no puede sumarse sin restricción a los 108 A máximos del rectificador. El control aplica $S_{serv}+S_{bancada}+S_{aux}\leq90$ kVA, con medición de todas las fases y descarga anticipada; la opción ALC de CT sobre servicio no es un limitador garantizado del total. Protecciones y selectividad se concilian con ELE-4.

\textbf{Reserva individual ampliada.} La obra nueva ofrece 2.000 L nominales y al menos 1.500 L útiles por grupo, reservando 200 L de expansión y hasta 300 L no extraíbles. A+B suma 3.000 L útiles sin compartirlos para acreditar un camino. Con consumo de referencia publicado de 41,3 L/h se obtienen $1.500/41,3=36,3196$ h, y la exigencia de veinticuatro horas con margen propio del veinte por ciento permite hasta $q_{lim}=1.500/(24\times1,20)=52,0833$ L/h. Este aumento elimina la dependencia del redondeo de 1.190 L, pero no constituye una ficha corregida de consumo: si la selección ambiental exige más, se redimensiona antes de compra. A $q_{lim}$ hay 28,8 h de reserva individual \parencite[pp.~1--3]{p7-fg-p200-ficha}.

El pedido se activa a 600 L útiles; con plazo declarado de ocho horas, a $q_{lim}$ quedan 183,3333 L o 3,52 h. A la referencia de 41,3 L/h quedan 269,6 L o 6,5278 h. Se conserva Enex Express como servicio declarado y Copec Empresas como respaldo de gestión, con cobertura, vigencia, acceso y descarga a formalizar antes de habilitación. No se acredita contrato ejecutado ni se descuenta su entrega para las primeras veinticuatro horas. El combustible corregido, potencia efectiva y respuesta a escalones de la instalación a 35\,$^{\circ}$C todavía requieren selección del fabricante; mayor estanque no demuestra esos parámetros.

\subsubsection{Balance conjunto y revisión del perfil anual (ECI-4)}
\label{sec:sd4-energia-eci4}

La incorporación de protección y monitorización se contabiliza en las mismas fronteras de salida y entrada. INC-4 fija techo propio de 0,650 kW/0,750 kVA para toda la electrónica, alarma, disparo, monitorización por PLC y recarga, en lugar de los 0,200 kW/0,2222 kVA anteriores; MON16 añade dos entradas máximas de 57,5 VA, sin adjudicarles FP unitario de catálogo. La cota anterior de 120 W se conserva exclusivamente para seguridad y acceso residual del recinto, excluyendo MON16; no se vuelve a contar su electrónica dentro de esos 120 W. Se adopta $P\leq S$ para estos monitores. Antes de añadir el tratamiento de bancos, el incremento de salida conservador es 0,565 kW/0,6428 kVA, o 0,678 kW/0,7713 kVA con el margen propio del veinte por ciento. La envolvente intermedia crece a 49,1696 kW/52,0048 kVA; su cociente conservador es 0,9455. Esta cifra no es el balance final de la instalación ambiental.

Se asigna incremento de incendio a L2 y 57,5 VA a cada una de L2/L3; no se recarga L1 con esas partidas. Las demandas anteriores de 19,3965/14,4140/17,4229 kVA ya incorporan el veinte por ciento. Al aplicar ese factor también al incremento, las nuevas cotas son 19,3965/15,1163/17,4919 kVA antes de tratamiento. El margen estacionario L1 hasta 20 kVA es sólo 0,6035 kVA; no equivale a reserva útil para cualquier ventilador, compresor o arranque. Cargas ambientales trifásicas o reparto nuevo requieren recomputar todas las fases y el caso de camino retirado; los 7,9952 kVA de margen total de placa no sustituyen esa comprobación.

La renovación de 900 m$^3$/h por banco de ENV16 incorpora 35,0322 kW entálpicos agregados, incluyendo 27,5752 kW latentes y 40,5187 kg/h de agua bajo las condiciones propias de cálculo. No se divide por COP 3 para atribuir un consumo a un equipo aún sin seleccionar. La electricidad de ventiladores, tratamiento, rehumidificación, controles y sus pérdidas integra salida protegida durante batería; conectarla antes de UPS crearía otra autonomía por resolver. El calor electroquímico y la disipación de las UPS en batería/recarga no se sustituyen por sus puntos AC/AC nominales. Si la combinación seleccionada excede potencia, fase o temperatura, se cambia UPS/banco y generación de manera conjunta, conservando la obligación de treinta minutos a plena carga.

El perfil anual se calcula por estados excluyentes red normal, batería, grupo, retorno con recarga y prueba; las horas suman 8.760 y los ciclos conservan reserva inicial/final equivalente. La energía de salida útil y la de recarga bruta no se suman como dos consumos del mismo servicio. Se registran por período energía de entrada, salida, DC hacia banco y desde banco, auxiliares interiores y exportación, además de baja carga de grupo y ensayos externos. Las pérdidas netas de banco se obtienen del ciclo, sin asignar los 35,28 kW DC enteros como calor real ni repetirlos dentro de la cota de 108 A.

Como actualización parcial del perfil anterior, la cota continua de incendio es 0,300 kW, frente a 0,200 kW presupuestados; monitorización aporta hasta 0,115 kW. El cambio neto interior normal es hasta 0,215 kW, equivalente a 1.883,40 kWh/año antes de pérdidas y clima. Para planificación propia se reservan 24 h/año acumuladas de alarma/recarga del panel a 0,650 kW; los 0,350 kW adicionales agregan 8,40 kWh, reemplazando la cota normal durante esas horas. Las pruebas de incendio usan su propia reserva y no se confunden con horas de corte de red. Ninguna de esas fracciones modifica máximos de compra. Las horas de tratamiento, ventilación y recarga de los bancos se incorporan una vez que exista una combinación seleccionada, con desempeño a carga baja y humedad; el perfil queda por cerrar.

Los valores 4,74 de referencia y 9,6715 de sensibilidad de GEN-4/ENE-4 permanecen identificados como estimaciones condicionadas anteriores: no son la PUE de la configuración dimensionada para 60 A y 1.800 m$^3$/h. ECI-4 conserva RT-06.11 por completar hasta declarar kW, FP y PUE de esa combinación completa. La medición prevista compara energía importada del recinto, descontando exportación y conservando las pérdidas internas, con TI interior en el mismo período. Los ensayos, las mediciones de autonomía y el funcionamiento ambiental se recibirán; no se presentan como ejecutados.
